\section{绪\hspace{1em}论}

\subsection{课题背景、目的与意义}

\subsubsection{研究背景}

大规模模型训练、容器化微服务、云端数据湖等场景有一个共同特点：产生并消费海量小文件。一次分布式训练作业可能在 checkpoint 阶段写出数十亿个参数分片，Kubernetes 集群的 Pod 在启动时也要创建和读取大量配置文件与临时卷。百度沧海存储团队在 Mantle\cite{li2025mantle} 中给出了一组数据：其云对象存储管理的文件对象数已达万亿量级，元数据请求规模每年仍有数倍增长。Ren 和 Gibson\cite{ren2013tablefs} 的测量显示，\texttt{stat}、\texttt{open}、\texttt{readdir} 等元数据操作占文件系统全部 I/O 操作的 60\%--80\%。在此类负载下，瓶颈通常不在数据传输带宽，而在元数据路径上：一旦热点目录下的文件数达到百万级，元数据查询与更新的延迟会直接压低系统吞吐上限。

\begin{figure}[htbp]
    \centering
    \begin{tikzpicture}[
        x=1cm, y=0.8cm,
        disk/.style={draw, fill=black!5, minimum width=1.5cm, minimum height=0.7cm,
                      align=center, font=\scriptsize},
        arr/.style={->, >=stealth, thick, black!50},
        lbl/.style={font=\footnotesize, text=black!60}
    ]

    % 入口
    \node[draw, rounded corners=2pt, fill=black!8, minimum width=4cm, minimum height=0.6cm,
          font=\small] (call) at (5.5, 6) {\texttt{create("/mnt/dir/newfile")}};

    % 6 个磁盘写入区域
    \node[disk] (jnl1) at (0.5, 3.8) {Journal\\[-1pt]描述符};
    \node[disk] (bmp)  at (2.8, 3.8) {inode\\[-1pt]bitmap};
    \node[disk] (itbl) at (5.1, 3.8) {inode 表\\[-1pt]（新文件）};
    \node[disk] (dblk) at (7.4, 3.8) {目录\\[-1pt]数据块};
    \node[disk] (pino) at (9.7, 3.8) {父目录\\[-1pt]inode};
    \node[disk] (jnl2) at (12, 3.8)  {Journal\\[-1pt]commit};

    % 步骤编号（带圈数字，放在磁盘框左侧以避开自顶向下的箭头）
    \node[font=\scriptsize\bfseries, anchor=east] at ([xshift=-2pt]jnl1.west) {\textcircled{\scriptsize 1}};
    \node[font=\scriptsize\bfseries, anchor=east] at ([xshift=-2pt]bmp.west)  {\textcircled{\scriptsize 2}};
    \node[font=\scriptsize\bfseries, anchor=east] at ([xshift=-2pt]itbl.west) {\textcircled{\scriptsize 3}};
    \node[font=\scriptsize\bfseries, anchor=east] at ([xshift=-2pt]dblk.west) {\textcircled{\scriptsize 4}};
    \node[font=\scriptsize\bfseries, anchor=east] at ([xshift=-2pt]pino.west) {\textcircled{\scriptsize 5}};
    \node[font=\scriptsize\bfseries, anchor=east] at ([xshift=-2pt]jnl2.west) {\textcircled{\scriptsize 6}};

    % 箭头
    \draw[arr] (call.south) -- (jnl1.north);
    \draw[arr] (call.south) -- (bmp.north);
    \draw[arr] (call.south) -- (itbl.north);
    \draw[arr] (call.south) -- (dblk.north);
    \draw[arr] (call.south) -- (pino.north);
    \draw[arr] (call.south) -- (jnl2.north);

    % 磁盘示意条
    \draw[fill=black!3] (0, 2.8) rectangle (12.5, 3.3);
    \node[lbl, anchor=east] at (-0.2, 3.05) {磁盘};

    % 写入位置标记
    \fill[black!25] (0.5, 2.85) rectangle (1.1, 3.25);
    \fill[black!25] (2.8, 2.85) rectangle (3.4, 3.25);
    \fill[black!25] (5.1, 2.85) rectangle (5.7, 3.25);
    \fill[black!25] (7.4, 2.85) rectangle (8.0, 3.25);
    \fill[black!25] (9.4, 2.85) rectangle (10.0, 3.25);
    \fill[black!25] (11.4, 2.85) rectangle (12.0, 3.25);

    % 跳转箭头
    \draw[<->, black!30] (1.1, 2.4) -- (2.8, 2.4);
    \draw[<->, black!30] (3.4, 2.4) -- (5.1, 2.4);
    \draw[<->, black!30] (5.7, 2.4) -- (7.4, 2.4);
    \draw[<->, black!30] (8.0, 2.4) -- (9.4, 2.4);
    \draw[<->, black!30] (10.0, 2.4) -- (11.4, 2.4);

    \node[lbl] at (6.25, 1.8) {$\geq 6$ 次随机 I/O，分布在互不相邻的磁盘区域};

    \end{tikzpicture}
    \caption{ext4 中一次 \texttt{create} 操作的磁盘写入路径}
    \label{fig:ext4-create-path}
\end{figure}

从文件系统的角度看，上述负载对 ext4、XFS 这类基于 B+ 树（或 ext4 的 H-tree 变体）的本地文件系统构成两个层面的压力。第一个层面是\textbf{粒度不匹配导致的写放大}。以 \texttt{create("/mnt/dir/newfile")} 为例，一次创建在 ext4 上至少触及五类磁盘区域：journal 描述符块、inode bitmap、inode 表中的新 inode 条目、父目录数据块中的新 dirent，以及父目录 inode 自身的时间戳与链接计数；事务结束时还要追加一条 commit 块。这五类区域分布在分区上互不相邻的位置，ordered 模式下日志区先写一遍、原位再写一遍，单次创建至少落 6 次随机 I/O（图~\ref{fig:ext4-create-path}）。一次创建真正需要写入的有效数据不超过 300 字节（256 字节 inode 加几十字节目录项），但 ext4 的写入粒度是 4\,KiB 磁盘页，即使只改 256 字节也要读出整页、修改后再写回。每秒数万次 \texttt{create} 时，写放大的累积效应会显著拖累系统吞吐。

第二个层面是\textbf{共享目录下的锁争用}。ext4 对目录数据块的修改受目录级互斥锁 \texttt{i\_rwsem} 保护，同一目录下的所有文件创建必须串行通过这把锁。底层即便是 NVMe SSD（Solid-State Drive）这类提供数十万 IOPS（Input/Output Operations Per Second）的设备，单目录下的创建吞吐也会被该锁压到远低于硬件能力的水平。上述两个层面共同决定了传统文件系统在元数据密集负载下的上限：粒度不匹配吃掉了单次操作的效率，目录锁则吃掉了并发度。

基于 LSM-tree（Log-Structured Merge-tree）\cite{oneil1996lsmtree} 的键值（KV）存储为该问题提供了一条不同的技术路线。KV 存储把操作粒度从 4\,KiB 磁盘页细化到单条记录：一个目录项即一个键值对，写入时只做一次内存追加，不涉及 bitmap 翻转、journal 双写或 4\,KiB 对齐的页填充；MemTable 写满之后一次性顺序刷盘，将原本分散的随机写归并为连续的顺序 I/O。TableFS\cite{ren2013tablefs} 在 2013 年把全部元数据存入 LevelDB\footnote{\url{https://github.com/google/leveldb}}，通过 FUSE（Filesystem in Userspace）挂载，小文件创建吞吐达到同期 ext4 的 2--3 倍。此后 RocksDB\cite{dong2021rocksdb} 在 LevelDB 基础上加入了列族、事务和前缀布隆过滤器等特性，Apache Ozone\footnote{\url{https://ozone.apache.org}} 和 Facebook Tectonic\cite{pan2021tectonic} 在 EB 级生产环境中把 RocksDB 用作元数据引擎，说明这条路径在工业规模下已经成立。不过上述生产系统都不对外暴露完整的 POSIX（Portable Operating System Interface）接口，键编码、父目录写放大、跨键事务等只在 POSIX 语义下才集中出现的问题，在它们的路径上并未完整展开。

\subsubsection{面临的问题和挑战}

将 KV 引擎用于文件系统元数据管理，需要解决三个核心设计问题（图~\ref{fig:three_challenges}）。

\textbf{问题一：树形目录到扁平键空间的映射。}
KV 引擎看到的是字节串到字节串的扁平映射，文件系统的命名空间却是一棵目录树。常见的目录访问可归为三种模式：\texttt{create} 与 \texttt{unlink} 改动一条目录项，对应 KV 的单键写或删；\texttt{lookup} 在给定父 inode 号和子文件名时要直接定位到一条记录，对应一次点查；\texttt{readdir} 要枚举某个目录下的全部子项，理想情况下同一父目录下的条目在键空间中物理相邻，只需一次前缀扫描。这三种模式对键编码提出的要求并不一致：以全路径为键对 \texttt{lookup} 最友好，但会让 \texttt{rename} 退化为整棵子树的级联重写；以父 inode 加文件名为键可以让 \texttt{rename} 保持常数代价，但只有在字节序与数值序一致的情况下，同一父目录下的条目才会在 LSM-tree 中连续，前缀扫描与前缀布隆过滤器也才能生效。如何同时兼顾这三种访问模式，是键编码必须回答的第一个问题。

\textbf{问题二：父目录属性更新形成的写热点。}
KV 存储把单次写入的粒度压到了一条记录，常规目录项的写放大由此缓解。但 POSIX 语义\cite{ieee2017posix} 还要求目录项变动时同步更新父目录的 \texttt{mtime} 与 \texttt{ctime}，\texttt{mkdir} 与 \texttt{rmdir} 还要改动 \texttt{nlink}。若沿用传统的读改写方式，每次 \texttt{create}、\texttt{unlink} 等操作都要先取出父目录 inode、修改几个字段、再写回同一条键。热点目录（如 \texttt{/tmp}、训练任务的 checkpoint 目录）下，这条键会承受大量并发的读改写：事务并发下冲突重试会累积，RocksDB 层面同一条 key 也会堆出大量历史版本使 compaction 反复归并，点查读放大同样上升。KV 写路径原本的收益在于把随机写整理为顺序追加，但读改写同一条键会把整个键空间的分散写收窄到这一条键上，父目录属性更新也就成了瓶颈。

\begin{figure}[htbp]
    \centering
    \begin{tikzpicture}[
        x=1cm, y=1cm,
        prob/.style={draw, rounded corners=2pt, fill=black!6, minimum width=3.6cm,
                      minimum height=0.7cm, align=center, font=\small},
        sol/.style={draw, rounded corners=2pt, fill=black!3, minimum width=3.6cm,
                     minimum height=0.7cm, align=center, font=\small},
        arr/.style={->, >=stealth, thick, black!50},
        lbl/.style={font=\scriptsize, text=black!55}
    ]

    % 列标题
    \node[font=\small\bfseries] at (1.8, 5.8) {设计问题};
    \node[font=\small\bfseries] at (10.2, 5.8) {本课题方案};

    % 第一行
    \node[prob] (p1) at (1.8, 4.8) {树 $\to$ 扁平键空间};
    \node[sol]  (s1) at (10.2, 4.8) {以边为中心 + 大端序};
    \draw[arr]  (p1) -- (s1);
    \node[lbl]  at (6, 5.15) {$O(1)$ lookup / readdir / rename};

    % 第二行
    \node[prob] (p2) at (1.8, 3.2) {父目录写放大};
    \node[sol]  (s2) at (10.2, 3.2) {DeltaOp 增量追加};
    \draw[arr]  (p2) -- (s2);
    \node[lbl]  at (6, 3.55) {写路径无需读取基值};

    % 第三行
    \node[prob] (p3) at (1.8, 1.6) {并发 TOCTOU 竞态};
    \node[sol]  (s3) at (10.2, 1.6) {悲观事务 + 行级锁};
    \draw[arr]  (p3) -- (s3);
    \node[lbl]  at (6, 1.95) {检查与操作在同一事务内};

    \end{tikzpicture}
    \caption{本课题面临的三个设计问题及对应解决方案}
    \label{fig:three_challenges}
\end{figure}

\textbf{问题三：并发目录操作的正确性。}
多个线程同时操作同一目录时，条件检查和执行修改之间存在一段时间窗口，即 TOCTOU（Time-of-Check-to-Time-of-Use）竞态。例如线程 A 执行 \texttt{rename} 时已确认目标路径不存在，在它写入新目录项之前线程 B 已在同一位置创建了同名文件；KV 层的写入是 last-write-wins 的，A 的操作会直接覆盖 B 的结果，外部观察者看到的语义不自洽。类似的窗口还出现在 \texttt{rmdir} 的空目录判断、\texttt{create} 的重名检查上。内核文件系统依赖 VFS 的 inode 互斥锁串行化这些操作，KV 引擎之上没有这层保护，必须由存储层事务把条件检查和修改动作收拢在同一视图下，否则单点正确的语义在并发下不再成立。

\subsubsection{课题目的与意义}

本课题的目标是设计并实现一个基于 RocksDB 的用户态 FUSE 文件系统 RucksFS，针对上述三个问题给出可落地的方案，并通过实验验证。元数据引擎选用 RocksDB，其 Column Family 可以把 inode 属性、目录项、增量记录放在相互独立的键空间中各自配置参数，前缀迭代器和前缀 Bloom 过滤器能够支撑目录前缀扫描，悲观事务（TransactionDB）提供了行级锁与冲突检测，恰好对应前述三个问题所需要的三类能力。对外接口选择 FUSE，主要理由是 FUSE 的回调与 POSIX 操作几乎一一对应，适合把元数据路径上的每一步实现清晰展示；与 JuiceFS\footnote{\url{https://juicefs.com}}、CubeFS\cite{liu2019cubefs} 等开源系统的接入方式一致，便于在相同工具链下进行性能对比。实现语言采用 Rust\cite{klabnik2023rust}：所有权与生命周期模型可以在编译期排除大部分内存安全问题，async/await 生态（tokio）也与 FUSE 的异步处理模式契合度较高。

围绕上述三个问题，本课题分别给出三项设计：键编码采用以边为中心的思路，将父 inode 号与子文件名按大端序拼接作为目录项键，使文件操作对应目录树中一条边的常数时间增删；父目录写放大问题通过增量更新 DeltaOp 解决，写入时只追加一条变化量不改动基础 inode 值，由后台线程定期将累积的增量折叠回基值；并发正确性由 RocksDB 悲观事务保证，将检查与修改封装在同一事务内，按 inode 编号排序加锁以预防死锁。三项设计的细节在第~\ref{sec:design}~章展开。

\textbf{课题意义。} 在研究意义上，RocksDB 自 6.0 版起提供了成熟的悲观事务接口，但把它系统性地应用到 POSIX 文件系统元数据层、并配上增量更新与合适键编码的公开工作仍较少；本课题为这一组合给出一个完整的设计与实现样本。在工程意义上，RucksFS 保持了轻量与可复现：整个系统跑在单机 RocksDB 之上，不依赖外部分布式事务数据库或跨服务的缓存一致性协议，便于后续在此骨架上增加元数据分片、副本容错或元数据预取等能力。课题的预期交付物包括一套可独立部署的 RucksFS 原型、覆盖 POSIX 合规性与元数据性能两条线的实验数据，以及用于复现实验的测试脚本与配置文件，全部公开在本课题配套仓库中。

\subsection{国内外研究现状}

用 KV 引擎管理文件系统元数据的研究始于 2013 年，至今已积累了从单机原型到 EB 级生产系统的实践。

\subsubsection{基于 KV 存储的元数据管理研究}

TableFS\cite{ren2013tablefs} 是这一方向最早的系统性工作。2013 年，Carnegie Mellon University 的 Ren 和 Gibson 把 ext4 的目录项和 inode 信息整体存入 LevelDB，通过 FUSE 挂载为标准文件系统。键以 “父 inode 号，文件名” 拼接构成，值为序列化后的 \texttt{struct stat} 加上可选的内联数据。在小文件场景下，TableFS 的 \texttt{create} 吞吐量达到同期 ext4 的 2--3 倍，但也暴露三处缺陷：\texttt{readdir} 需要扫描 LevelDB 的全部键空间，文件总数百万级时延迟明显上升；LevelDB 不提供事务语义，多线程写入的正确性完全依赖外部加锁；inode 号按原生字节序存储，同一目录下的条目在 LSM-tree 中不一定连续，前缀迭代器与前缀布隆过滤器无法生效。TableFS 也没有处理父目录 inode 更新带来的写放大问题。这三处缺陷分别对应本课题在键编码、事务和写放大三个方向的改进。

同一团队次年推出了 IndexFS\cite{ren2014indexfs}，把 TableFS 思路扩展到分布式环境。每个元数据节点上运行一个 LevelDB 实例，目录项通过 GIGA+ 动态哈希分区算法分布到不同节点。在 SST 文件的内部组织上，IndexFS 采用列式（column-style）布局：将频繁访问的查找属性与完整 inode 元数据分别放在 Index 表和 Log 表中，Index 表参与 Compaction，Log 表则保留追加写形式不参与，从而降低维护开销。借助这一布局，IndexFS 在 128 个元数据服务器上达到约 84.2 万次/秒的 \texttt{create} 吞吐，论文获 SC'14 最佳论文奖。LocoFS\cite{li2017locofs} 在 2017 年针对目录树紧耦合导致的访问效率问题，提出松耦合架构：单个目录元数据服务（DMS）维护 \texttt{path → d-inode} 映射，多个文件元数据服务（FMS）按 \texttt{(dir\_uuid + file\_name)} 一致性哈希分布文件 inode；文件 inode 进一步拆分为访问属性与内容属性两部分，以减小单次操作的粒度。Xu 等人\cite{xu2014drop} 在 2014 年提出了面向 EB 级规模的子树粒度元数据划分方案。这些工作侧重多节点元数据划分与扩展策略，与本课题的关注点（在 RocksDB 之上如何组织 POSIX 元数据操作）不在同一层面。

百度沧海·存储团队在 SOSP'25 上发表的 Mantle\cite{li2025mantle} 是这一路线的最新进展。Mantle 为云对象存储提供层级命名空间支持，其底层 KV 引擎 TafDB 上设计了两类元数据表：MetaStore 以 \texttt{(parent\_id, name)} 为主键并以 \texttt{inode\_id} 建立二级索引，分别支撑 \texttt{lookup}/\texttt{readdir} 与按 inode 的 \texttt{getattr}；IndexNode 表则以全路径作为键，配合内存中的 TopDirPathCache 加速顶层稳定路径解析。在并发更新方面，Mantle 引入 Delta Record 机制：当多个客户端在同一目录下并发执行 \texttt{mkdir}/\texttt{rmdir} 等修改父目录属性的操作时，不再以读改写方式更新原记录，而是各自追加一条以 “原属性键加时间戳” 为键的增量记录，由后台线程异步合并到基准值，从而消除热点目录属性上的写竞争。Mantle 依赖的 TafDB 是百度自研的类 Spanner 分布式事务数据库，IndexNode 与 TafDB 之间还要维持一套较复杂的缓存一致性协议，整套系统并不开源。本课题借鉴 Mantle 的增量更新思路，把它实现在 MetadataServer 本地的 RocksDB 事务上：DeltaOp 追加、懒折叠和折叠结果缓存全部在同一 RocksDB 实例内完成，不再依赖外部分布式事务数据库和跨服务的缓存一致性协议，以较低的系统复杂度复现核心收益。

\subsubsection{云原生文件系统中的元数据方案}

早期的分布式文件系统以集中式元数据管理为主：GFS\cite{ghemawat2003gfs} 和 HDFS\cite{shvachko2010hdfs} 把全部元数据驻留在单节点内存中，Ceph\cite{weil2006ceph} 用动态子树划分将目录树分配到多个元数据服务器。这些系统在 PB 级规模上运转良好，但元数据可扩展性始终受限于内存容量或子树迁移开销。分布式 KV 存储（以 Amazon Dynamo\cite{decandia2007dynamo} 的可用性优先设计为代表）逐步成为上层文件系统的新型后端选择之后，Ozone 与 Tectonic 在 2020 年前后将元数据引擎替换为 RocksDB，分别管理 EB 级数据，但都不提供完整 POSIX 接口。

JuiceFS\footnote{\url{https://juicefs.com}} 是国内使用较广的开源 POSIX 文件系统，元数据存储在 Redis、TiKV 或关系型数据库（MySQL、PostgreSQL 等）之中，数据块写入 S3 兼容的对象存储。这种解耦方式允许用户按需选择后端组合，代价是性能与可用性完全取决于外部组件：Redis 单线程模型在元数据操作超过每秒数万次时会出现瓶颈，TiKV 则引入 Raft 共识协议带来的额外延迟。JuiceFS 的 FUSE 客户端使用 Go 实现，运行时垃圾回收在极端负载下可能引入延迟抖动。CubeFS\cite{liu2019cubefs}（最初名为 CFS）于 2019 年在 SIGMOD 发表，后捐赠给 CNCF。与 JuiceFS 不同，CubeFS 自带元数据引擎，每个元数据分区使用内存中的 B-tree 维护目录结构，多副本之间通过 Multi-Raft 协议保持一致性。CubeFS 的读路径延迟在微秒级，写路径受 Raft 日志同步所限典型在毫秒级；内存 B-tree 也限制了单节点可管理的元数据量，按每个 inode 约 300 字节估算，十亿文件需要约 280\,GB 内存。

上述系统按元数据承载方式可归为三类：内存驻留型（GFS、HDFS、CubeFS）、外部数据库型（JuiceFS）、嵌入式 KV 型（Ozone、Tectonic）。三类方案侧重点不同，但都把分布式部署作为前提，要么把一致性和事务交给外部数据库，要么把复制、分片和故障恢复做进文件系统自身。一个更具体的问题随之出现：在已经采用分布式部署的前提下，元数据引擎本身该如何围绕文件操作做定制，而不是完全受制于通用数据库或共识框架的语义与开销。

\subsubsection{元数据引擎的深度优化研究}

与分片扩展路线并行的另一条研究路线更关注元数据引擎本身的执行效率。SingularFS\cite{guo2023singularfs} 在 ATC'23 上展示了一个激进的设计：仅用单个元数据服务器支撑十亿级文件规模的分布式文件系统。其核心优化包括无日志元数据操作、层次化并发控制以及 NUMA（Non-Uniform Memory Access）感知的 inode 分区。在 RDMA（Remote Direct Memory Access）网络和持久内存的硬件环境下，SingularFS 的单节点元数据吞吐量远超传统多节点方案，说明单机元数据服务本身仍有较大的性能潜力。本课题的 MetadataServer 同样采用单实例部署，但不依赖 RDMA 与持久内存，而是以 RocksDB 作为载体，借助其 WAL（Write-Ahead Logging）和事务保证崩溃一致性，代价是单次操作延迟会高一些。Wang 等人\cite{wang2023cfs} 在 EuroSys'23 上发表的 CFS 通过数据布局优化与单分片原子原语，把跨节点事务收缩为单分片内的原子操作，在 50 节点集群上相对 HopsFS 等先前方案的整体吞吐提升达到一个数量级以上。本课题在锁粒度上选了类似的方向：使用 RocksDB 悲观并发控制（Pessimistic Concurrency Control, PCC）事务自带的行级锁，只对实际访问到的 inode 行或目录项键加锁，而不是整个父目录。

Yang 等人\cite{yang2020batchfile} 提出的批量文件操作 BFO 把多个小文件的读取与写入合并执行，以摊薄系统调用与磁盘 I/O 的固定开销。本课题的 \texttt{readdirplus} 沿用 FUSE 协议本身一次返回目录项加属性的形式，没有进一步在客户端侧实现跨多文件的批量合并。FetchBPF\cite{mishler2024fetchbpf} 通过扩展 eBPF 框架在 Linux 内核中支持可定制的内存预取策略，FBMM\cite{tabatabai2024fbmm} 则把内存管理器实现为文件系统，利用 VFS 的可扩展性引入新策略。这两项工作与本课题的元数据组织并不直接相关，但说明 eBPF 与 VFS 都是后续在元数据预取或扩展非标准语义时的可行通道。MetaHive\cite{heidari2024metahive} 针对通用 KV 缓存替换策略未考虑文件系统访问局部性的问题设计了元数据感知的缓存分区方案；FUSEE\cite{shen2023fusee}、Duplex\cite{dong2023duplex} 与 SwitchFS\cite{xu2026switchfs} 分别在内存解耦架构、全路径索引和交换机侧异步协调上给出新的尝试。这几项工作依赖的硬件和体系结构与本课题的 RocksDB 加 FUSE 组合相距较远，但 FUSEE 把元数据拆成可远程访问的数据结构这一做法，对本课题后续做元数据分片仍有参考意义。

国内方面，范立衡等人\cite{fan2014kvmetadata} 较早探索了 KV 存储用于元数据集群副本一致性的管理策略，其对哪些字段值得副本间强一致、哪些可以放宽的分类，对 RucksFS 未来引入副本时有参考价值。吴雨飞等人\cite{wu2025dfscache} 在 2025 年针对基于 TiKV 的分布式文件系统元数据缓存一致性问题，提出了并发广播与惰性删除相结合的轻量级维护机制，可直接借用到 RucksFS 的客户端缓存放宽方向。

\subsubsection{现有研究与本课题的定位对比}

表~\ref{tab:system_compare} 从 KV 引擎、事务支持、增量更新和用户态接口四个维度，把本课题与最相关的几个系统并列对比。

\begin{table}[!ht]
    \centering
    \caption{本课题与相关系统的定位对比}
    \label{tab:system_compare}
    \small
    \begin{tabular}{lcccc}
    \toprule
        系统 & KV 引擎 & 事务支持 & 增量更新 & 用户态接口 \\ \midrule
        TableFS\cite{ren2013tablefs} & LevelDB & 无 & 无 & FUSE \\
        IndexFS\cite{ren2014indexfs} & LevelDB & WriteBatch & 无 & 自定义库 \\
        Mantle\cite{li2025mantle} & TafDB & 分布式事务 & Delta Record & 对象存储 API \\
        JuiceFS & 多种后端 & 依赖后端 & 无 & FUSE \\
        CubeFS\cite{liu2019cubefs} & 内存 B-tree & Raft & 无 & FUSE \\
        RucksFS（本课题） & RocksDB & PCC 事务 & DeltaOp & FUSE \\
    \bottomrule
    \end{tabular}
\end{table}

\subsection{论文的主要内容与结构}

\subsubsection{主要工作}

本课题围绕基于 KV 存储的文件系统元数据管理这一目标，主要完成以下五项工作。

（1）设计并实现基于 RocksDB 多列族的元数据存储方案。inode 属性、目录项、增量记录、系统元数据分别放在独立列族中，各自配置压缩策略和缓存参数。目录项列族的键采用以边为中心的编码，父 inode 号和子文件名按大端序拼接，每条目录项对应目录树中的一条父子边；\texttt{create}、\texttt{unlink}、\texttt{rename} 均为常数时间。大端序使字节序与数值序一致，同一父目录下的子项在键空间中相邻，\texttt{readdir} 通过前缀扫描完成，\texttt{lookup} 通过点查完成。

（2）引入增量更新与懒折叠（lazy folding）机制，并以 RocksDB 悲观事务统一处理元数据并发。文件创建和删除不直接修改父目录 inode，而是以 DeltaOp 形式追加一条增量记录；读路径在 \texttt{getattr} 或 \texttt{readdir} 时将基值与增量在线合并；后台线程在增量数量超过 32 条时折叠回基值。所有元数据变更操作通过 RocksDB 悲观事务完成加锁与冲突检测，事务冲突时服务端自动重试最多 3 次，仍失败则在 FUSE 层映射为 \texttt{EAGAIN} 返回。

（3）基于 FUSE 实现覆盖 \texttt{create}、\texttt{mkdir}、\texttt{unlink}、\texttt{rmdir}、\texttt{rename}、\texttt{readdir}、\texttt{setattr}、\texttt{read}、\texttt{write} 等操作的完整接口集合。客户端基于 Rust 的 fuse3 库与 tokio 异步运行时实现，挂载时启用 \texttt{default\_permissions} 与 \texttt{allow\_other}，常规权限检查交由内核 VFS 处理。

（4）实现由 Client、MetadataServer、DataServer 三段组成的分布式部署形态，三者通过 gRPC（tonic）通信。客户端内部设有请求分发层，把元数据请求交给 MetadataServer、数据请求交给 DataServer；\texttt{open}、\texttt{release} 这类跨服务的语义由客户端统一编排。

（5）完成正确性与性能两条线的系统评估。正确性方面，在分布式部署下运行 pjdfstest 套件共 398 条用例，按本课题实现范围分解后在实现范围内的用例上全部通过；性能方面，在 1 台 64 核服务器加最多 96 台 2 核客户端的集群上使用 mdtest 压测，中低并发区间（$N \in \{2, 8, 32\}$）将 RucksFS 与 JuiceFS+TiKV、NFS（Network File System）v4.2 作三方横向对比，高并发区间（$N \in \{64, 96\}$）改用 Delta 与 NoDelta 两个构建变体作内部对照。

实现的完整源代码、测试编排脚本、原始测量数据以及本论文的 LaTeX 源码均公开在 \url{https://github.com/csjgg/rucksfs}，按第~\ref{sec:experiments}~章的参数即可复现。

\subsubsection{论文结构}

本论文共六章。第一章绪论交代研究背景，分析 KV 存储用于文件系统元数据时需要面对的技术挑战，并梳理国内外相关工作。第二章相关技术基础介绍 FUSE 工作原理、RocksDB 的 LSM-tree 架构与列族/事务机制，以及 POSIX 文件系统语义的关键约束。第三章元数据组织与事务模型展开 MetadataServer 内部的设计：多列族的命名空间划分、目录项的大端序边键编码、DeltaOp 增量追加与懒折叠机制、PCC 事务封装与冲突重试策略，以及各 POSIX 目录操作在事务层面的具体落实。第四章系统实现把第三章的设计实现为一个完整的分布式原型，包括客户端的请求分发、跨服务操作编排、MetadataServer 与 DataServer 各自的实现，以及 gRPC 多通道连接池等优化。第五章测试与结果分析在 pjdfstest 与 mdtest 两条线上检验前文设计。第六章总结与展望回顾主要成果，讨论当前设计的局限和后续可扩展的改进方向。
